\documentclass[../../main.tex]{subfiles}

\begin{document}

\section{$X \times Y$ 上的积拓扑}\label{section:ch02-积拓扑XY}

若 $X$ 和 $Y$ 是两个拓扑空间，则有一个在笛卡儿积 $X \times Y$ 上定义拓扑的标准方法。下面就来研究这个拓扑及它的一些性质。

\begin{definition}[积拓扑]\label{definition:ch02-积拓扑XY}
设 $X$ 和 $Y$ 是两个拓扑空间，$X \times Y$ 上的积拓扑（product topology）是以族 $\mathcal{B}$ 为基的拓扑，其中 $\mathcal{B}$ 是所有形如 $U \times V$ 的集合的族，$U$ 和 $V$ 分别是 $X$ 和 $Y$ 的开子集。
\end{definition}

我们来验证 $\mathcal{B}$ 是一个基。第一个条件是显然的，这是由于 $X \times Y$ 本身就是一个基元素。第二个条件也是易于验证的。事实上，任意两个基元素 $U_{1} \times V_{1}$ 与 $U_{2} \times V_{2}$ 的交是
\begin{align*}
(U_{1} \times V_{1}) \cap (U_{2} \times V_{2}) = (U_{1} \cap U_{2}) \times (V_{1} \cap V_{2}),
\end{align*}
由于 $U_{1} \cap U_{2}$ 和 $V_{1} \cap V_{2}$ 分别是 $X$ 和 $Y$ 的开集，所以上述集合是一个基元素，参见\cref{figure:ch02-积基交}。

\par
\begin{figure}[htbp]
\centering
\includegraphics[width=0.45\textwidth]{ch02-fig15-1}
\caption{积基元素的交}\label{figure:ch02-积基交}
\end{figure}

值得注意的是：族 $\mathcal{B}$ 不是 $X \times Y$ 的一个拓扑。例如，\cref{figure:ch02-积基交}中两个矩形的并就不是两个集合的积，因而不属于 $\mathcal{B}$，但它是 $X \times Y$ 中的开集。

我们每次引进一个新概念，总是试图弄清它与前面引进的概念之间的关联。我们现在要问：当 $X$ 和 $Y$ 的拓扑是由它们的基给出时，关于积拓扑能说些什么呢？下面的定理给出了回答：

\begin{theorem}[积拓扑的基]\label{theorem:ch02-积基}
若 $\mathcal{B}$ 是 $X$ 的拓扑的一个基，$\mathcal{C}$ 是 $Y$ 的拓扑的一个基，则族
\begin{align*}
\mathcal{D} = \{B \times C \mid B \in \mathcal{B} \text{ 并且 } C \in \mathcal{C}\}
\end{align*}
是 $X \times Y$ 的拓扑的一个基。
\end{theorem}

\begin{proof}
我们应用\cref{lemma:ch02-开集族为基}。给定 $X \times Y$ 的一个开集 $W$ 以及 $W$ 的一个点 $x \times y$。根据积拓扑的定义，存在一个基元素 $U \times V$，使得 $x \times y \in U \times V \subset W$。因为 $\mathcal{B}$ 和 $\mathcal{C}$ 分别是 $X$ 和 $Y$ 的基，所以我们可以在 $\mathcal{B}$ 中选取一个元素 $B$，使得 $x \in B \subset U$，也可以在 $\mathcal{C}$ 中选取一个元素 $C$，使得 $y \in C \subset V$。于是 $x \times y \in B \times C \subset W$。从而族 $\mathcal{C}$ 符合\cref{lemma:ch02-开集族为基}的条件，因此 $\mathcal{D}$ 是 $X \times Y$ 的一个基。
\end{proof}

\begin{example}\label{example:ch02-R标准积}
我们有 $\mathbb{R}$ 的一个标准拓扑——序拓扑。这个拓扑与它自身的积称之为 $\mathbb{R} \times \mathbb{R} = \mathbb{R}^{2}$ 上的标准拓扑。它有一个基是 $\mathbb{R}$ 上所有开集的积的族，而刚才证明的\cref{theorem:ch02-积基}告诉我们：$\mathbb{R}$ 中所有开区间的积 $(a, b) \times (c, d)$ 这个更小的族也可以作为 $\mathbb{R}^{2}$ 的拓扑的一个基。每一个这样的积可以画成 $\mathbb{R}^{2}$ 中一个矩形的内部。因此 $\mathbb{R}^{2}$ 的标准拓扑正好是第\cref{section:ch02-拓扑的基}节例 2 中研究过的那一种。
\end{example}

有时也需要用子基来表示积拓扑。为此，我们先定义某些叫做投射的函数。

\begin{definition}[投射]\label{definition:ch02-投射}
设 $\pi_{1}: X \times Y \rightarrow X$ 定义为
\begin{align*}
\pi_{1}(x, y) = x.
\end{align*}
$\pi_{2}: X \times Y \rightarrow Y$ 定义为
\begin{align*}
\pi_{2}(x, y) = y.
\end{align*}
映射 $\pi_{1}$ 和 $\pi_{2}$ 分别称为 $X \times Y$ 到它的第一因子和第二个因子上的投射（projections）。
\end{definition}

我们使用“到……上”这个词，是因为 $\pi_{1}$ 和 $\pi_{2}$ 都是满射（除去 $X$ 和 $Y$ 中有一个是空集的情形，因为这时 $X \times Y$ 是空集，整个讨论就没有意义了！）。设 $U$ 是 $X$ 的一个开子集，那么集合 $\pi_{1}^{-1}(U)$ 恰好就是 $U \times Y$，它是 $X \times Y$ 中的开集。同样地，若 $V$ 是 $Y$ 的一个开集，则
\begin{align*}
\pi_{2}^{-1}(V) = X \times V
\end{align*}
也是 $X \times Y$ 中的开集。这两个集合的交是 $U \times V$，参见\cref{figure:ch02-投射原像}。这个事实引导出下面的定理：

\par
\begin{figure}[htbp]
\centering
\includegraphics[width=0.55\textwidth]{ch02-fig15-2}
\caption{投射的原像}\label{figure:ch02-投射原像}
\end{figure}

\begin{theorem}[积拓扑的子基]\label{theorem:ch02-积子基}
族
\begin{align*}
\mathcal{S} = \{\pi_{1}^{-1}(U) \mid U \text{ 是 } X \text{ 中的开集}\} \cup \{\pi_{2}^{-1}(V) \mid V \text{ 是 } Y \text{ 中的开集}\}
\end{align*}
是 $X \times Y$ 的积拓扑的一个子基。
\end{theorem}

\begin{proof}
用 $\mathcal{T}$ 表示 $X \times Y$ 的积拓扑，设 $\mathcal{T}'$ 是由 $\mathcal{S}$ 生成的拓扑，因为 $\mathcal{S}$ 的每一个元素都属于 $\mathcal{T}$，所以 $\mathcal{S}$ 的元素的有限交的任意并也属于 $\mathcal{T}$。因此，$\mathcal{T}' \subset \mathcal{T}$。另一方面，拓扑 $\mathcal{T}$ 的任意基元素 $U \times V$ 是 $\mathcal{S}$ 中元素的有限交，这是因为
\begin{align*}
U \times V = \pi_{1}^{-1}(U) \cap \pi_{2}^{-1}(V),
\end{align*}
所以 $U \times V$ 属于 $\mathcal{T}'$，因此 $\mathcal{T} \subset \mathcal{T}'$。
\end{proof}

\end{document}